\documentclass{article}
\usepackage{amsmath}
\usepackage{amssymb}
\usepackage{listings}
\usepackage{amsfonts} % For \mathcal
\usepackage[utf8]{inputenc} % For Slovenian characters

\lstset{
	language=Rust,
	escapeinside={(*}{*)},
	numbers=left,
	basicstyle=\ttfamily
}

\begin{document}

\section*{Referenca za relacije -- originalna formulacija}

\subsection*{1. Osnovne množice in elementi}
\begin{itemize}
	\item $\mathcal{L}$: Množica vseh možnih posoj (loans). Posamezno posojo označimo z $L$.
	\item $\mathcal{R}$: Množica vseh regij (regions) (spremenljivke regij, npr. '0, '1, ...). Posamezno regijo označimo z $R$. Regija predstavlja množico posoj. Torej velja $\mathcal{R} \subset 2^{\mathcal{L}}$
	\item $\mathcal{S}$: Množica vseh stavkov (statements) v MIR. Stavek $stmt \in \mathcal{S}$ je identificiran kot $BB_{i,j}$, $j$-ti stavek v $i$-tem osnovnem bloku (basic block).
	\item $\mathcal{P}$: Množica vseh točk (points) v grafu poteka kontrol (control flow graph - CFG) programa. Tocke so lahko na zacetku ali sredi stavka.
	      \begin{itemize}
		      \item $S(stmt)$: Točka na začetku stavka $stmt$ (start of statement) (pred izvedbo).
		      \item $M(stmt)$: Točka na sredini stavka $stmt$ (middle of statement) (med izvedbo).
	      \end{itemize}
	\item $C = (C_V, C_E)$: Graf poteka kontrol (Control Flow Graph), kjer je $C_V = \mathcal{P}$ množica točk (vozlišč) in $C_E \subseteq \mathcal{P} \times \mathcal{P}$ množica povezav poteka kontrol. Dodatno:
	      \begin{itemize}
		      \item Za vsak $stmt \in \mathcal{S}$, velja $(S(stmt), M(stmt)) \in C_E$.
		      \item Povezave obstajajo od $M(stmt)$ do začetnih točk $S(stmt')$ naslednjih stavkov $stmt'$ v CFG.
	      \end{itemize}
\end{itemize}

\subsection*{2. Začetne relacije (Input Relations)}
Te relacije predstavljajo začetni vhod za analizo. Izpeljana so iz sistema tipov in prejsnje analize v MIR.
\begin{itemize}
	\item \textbf{Začetna relacija vsebovanosti (Base Subset)} ($\beta \subseteq \mathcal{R} \times \mathcal{R} \times \mathcal{P}$):
	      \begin{itemize}
		      \item $(R_1, R_2, P) \in \beta$ pomeni, da sistem tipov zahteva, da je regija $R_1$ podmnožica regije $R_2$ (označeno $R_1 \subseteq R_2$ ali $R_1 : R_2$) na točki $P$. To dejstvo je vzpostavljeno na sredini točke ($P = M(stmt)$) stavka $stmt$, ki inducira zahtevo.
	      \end{itemize}

	\item \textbf{Začetna relacija posoje regij (Borrow Region)} ($\psi \subseteq \mathcal{R} \times \mathcal{L} \times \mathcal{P}$):
	      \begin{itemize}
		      \item $(R, L, P) \in \psi$ pomeni, da izraz za posojo (borrow expression) na točki $P$ ustvari posojo $L$ in jo poveže z regijo $R$. To dejstvo je vzpostavljeno na sredini stavka ($P = M(stmt)$).
	      \end{itemize}

	\item \textbf{Relacija žive regije (Region Live At)} ($\nabla_{R} \subseteq \mathcal{R} \times \mathcal{P}$):
	      \begin{itemize}
		      \item $(R, P) \in \nabla_{R}$ pomeni, da je regija $R$ živa (live) na točki $P$ (tj. spremenljivka, katere tip vključuje $R$, bo morda kasneje uporabljena). To določi standardna analiza živosti (liveness analysis) (podrobnosti tukaj izpuščene, glej NLL RFC).
	      \end{itemize}

	\item \textbf{Relacija prekinitve posoje (Killed Loan)} ($\kappa \subseteq \mathcal{L} \times \mathcal{P}$):
	      \begin{itemize}
		      \item $(L, P) \in \kappa$ pomeni, da je posoja $L$ "prekinjena" (killed) na točki $P$. To se običajno zgodi na sredini točke ($P = M(stmt)$) prireditvenega stavka $stmt$, ki prepiše pot (path), prej povezano s posojo $L$.
	      \end{itemize}

	\item \textbf{Relacija razveljavitve posoje (Invalidates Loan)} ($\iota \subseteq \mathcal{P} \times \mathcal{L}$):
	      \begin{itemize}
		      \item $(P, L) \in \iota$ pomeni, da dejanje, ki se zgodi na točki $P$ (npr. mutacija), razveljavi pogoje posoje $L$.
	      \end{itemize}
\end{itemize}

\subsection*{3. Izpeljane relacije (Derived Relations) (Izračunane z analizo)}

\begin{itemize}
	\item \textbf{Relacija vsebovanosti (Subset Relation)} ($\Gamma \subseteq \mathcal{R} \times \mathcal{R} \times \mathcal{P}$):
	      \begin{itemize}
		      \item $(R_1, R_2, P) \in \Gamma$ pomeni, da analiza sklepa, da mora biti $R_1$ podmnožica $R_2$ na točki $P$. Definirana je z zaprtjem naslednjih pravil:
		      \item \textbf{Osnova:} Če $(R_1, R_2, P) \in \beta$, potem $(R_1, R_2, P) \in \Gamma$.
		      \item \textbf{Tranzitivnost:} Če $(R_1, R_2, P) \in \Gamma$ in $(R_2, R_3, P) \in \Gamma$, potem $(R_1, R_3, P) \in \Gamma$.
		      \item \textbf{Propagacija (ze z dodatnimi pravili):} Če $(R_1, R_2, P) \in \Gamma$, $(P, Q) \in C_E$, $(R_1, Q) \in \nabla_R$, in $(R_2, Q) \in \nabla_R$, potem $(R_1, R_2, Q) \in \Gamma$. (Propagira se samo, če sta obe regiji živi na naslednji točki $Q$).
	      \end{itemize}

	\item \textbf{Relacija zahteve (Requires Relation)} ($\zeta \subseteq \mathcal{R} \times \mathcal{L} \times \mathcal{P}$):
	      \begin{itemize}
		      \item $(R, L, P) \in \zeta$ pomeni, da regija $R$ zahteva, da pogoji posoje $L$ veljajo na točki $P$. Definirana je z zaprtjem naslednjih pravil:
		      \item \textbf{Osnova:} Če $(R, L, P) \in \psi$, potem $(R, L, P) \in \zeta$.
		      \item \textbf{Vsebovanost:} Če $(R_1, L, P) \in \zeta$ in $(R_1, R_2, P) \in \Gamma$, potem $(R_2, L, P) \in \zeta$. (Če $R_1 \subseteq R_2$, potem $R_2$ zahteva vse, kar zahteva $R_1$).
		      \item \textbf{Propagacija:} Če $(R, L, P) \in \zeta$, $(L, P) \notin \kappa$, $(P, Q) \in C_E$, in $(R, Q) \in \nabla_R$, potem $(R, L, Q) \in \zeta$. (Propagira se samo, če posoja $L$ ni prekinjena na $P$ in je regija živa na naslednji točki $Q$).
	      \end{itemize}

	\item \textbf{Relacija živosti posoje (Loan Live At Relation)} ($\nabla_{L} \subseteq \mathcal{L} \times \mathcal{P}$):
	      \begin{itemize}
		      \item $(L, P) \in \nabla_{L}$ pomeni, da je posoja $L$ živa (live) na točki $P$.
		      \item \textbf{Definicija:} $(L, P) \in \nabla_{L}$ če $\exists R \in \mathcal{R}$ tako da $(R, P) \in \nabla_{R}$ in $(R, L, P) \in \zeta$. (Posoja je živa, če jo zahteva neka živa regija).
	      \end{itemize}
\end{itemize}

\subsection*{4. Pogoji za napako}

\begin{itemize}
	\item \textbf{Relacija napake (Error Relation)} ($\epsilon \subseteq \mathcal{P}$):
	      \begin{itemize}
		      \item $P \in \epsilon$ pomeni, da se napaka preverjevalnika izposoj (borrow check error) pojavi na točki $P$.
		      \item \textbf{Definicija:} $P \in \epsilon$ če $\exists L \in \mathcal{L}$ tako da $(P, L) \in \iota$ in $(L, P) \in \nabla_{L}$. (Napaka se pojavi, če dejanje na točki $P$ razveljavi posojo $L$, ki je živa na točki $P$).
	      \end{itemize}
\end{itemize}

\end{document}
